\documentclass[11pt]{article}
\usepackage[margin=1in]{geometry}
\usepackage{amsmath,amssymb,amsthm}
\usepackage[hidelinks]{hyperref}
\newtheorem{lemma}{Lemma}
\newtheorem{theorem}[lemma]{Theorem}
\title{A second thin-digit terminal family for ordinary gap six}
\author{Internal proof draft; no priority or external-review claim}
\date{18 November 2026}
\begin{document}\setlength{\emergencystretch}{1em}\maketitle
Fix an integer $r\ge1$ and put $L=7^r$, $N=2L+1$, $c=N+5$.
Here $N$ is the terminal order.
Work in a complete algebraically closed extension $K$ of
$\mathbb Q_7$, with valuation ring $\mathcal O$ and residue
field $k=\mathcal O/\mathfrak m$.
Let $A=c-zD_z+z$ and
$\Pi_N(X)=\prod_{\ell=0}^{N-1}(X+6\ell)$.
For $Q\in\mathcal O[X]$ of degree at most four define
\[
 G_Q(X)=\binom{X/6+N-1}{N}Q(X),\quad
 d_j(Q)=\sum_{\ell=0}^j(-1)^\ell\binom j\ell G_Q(c-\ell).
\]
For every integer $\ell$, the binomial argument
$(c-\ell)/6+N-1$ lies in $\mathbb Z_7$.
Binomial polynomials are integer-valued there, by continuity
from nonnegative integers; hence $d_j(Q)\in\mathcal O$.
The six normalized boundary rows, at $z^6=7v$, are
\[
 V_a(Q;v)=\sum_{\ell\ge0}
       \frac{7^\ell a!}{(a+6\ell)!}d_{a+6\ell}(Q)v^\ell,
                     \qquad0\le a<6.                    \tag{1}
\]
They lie in $\mathcal O[v]$; the five basis columns
$Q=1,X,\ldots,X^4$ lie in $\mathbb Z_7[v]$.
Indeed the finite formal identity
$\Pi_N(A)Q(A)1/(6^NN!)=\sum_jd_j(Q)z^j/j!$ holds, and
the coefficient weight in (1) has valuation
$(s_7(a+6\ell)-a)/6\ge0$.
Thus their simultaneous zero is exactly the original cyclic rank
relation. The argument uses formal finite differences and never
evaluates an exponential at a7-adic boundary point.

\begin{lemma}[The exact diagonal series]\label{lem:diagonal}
Let $R(t)$ be the unique series with $R(0)=1$ and $R-R^7=t$.
For the auxiliary polynomials
\[
 B_{0,\gamma,6}=1,\quad
 B_{j+1,\gamma,6}=(z+6j+\gamma+5)B_{j,\gamma,6}
                                      -zB'_{j,\gamma,6},
\]
one has, formally over $\mathbb Q[\eta,z][[t]]$,
\[
 \sum_{n\ge0}\frac{B_{n,n+\eta,6}(z)}{6^n n!}t^n
 =\frac{6R^{1-\eta}}{7R^6-1}\exp(z(1-R)).                \tag{2}
\]
For $\eta=0$, $R$ and $\kappa=6R/(7R^6-1)$ are7-adically integral,
and $\kappa\equiv R-7t\pmod{49}$.
\end{lemma}
\begin{proof}
The fixed-shape exponential generating series is
\[
 \sum_{n\ge0}\frac{B_{n,\gamma,6}(z)}{6^nn!}s^n
 =(1-s)^{-(\gamma+5)/6}
                \exp\bigl(z(1-(1-s)^{1/6})\bigr).
\]
Apply the formal diagonal Lagrange identity
\[
 \sum_{n\ge0}t^n[s^n]\alpha(s)^n\beta(s)
       =\frac{\beta(w)}{1-w\alpha'(w)/\alpha(w)},
       \qquad w=t\alpha(w),
\]
with $\alpha=(1-s)^{-1/6}$.
This identity follows directly by the formal residue substitution
$t=w/\alpha(w)$.
Putting $R=(1-w)^{1/6}$ gives $t=R-R^7$ and (2).
The implicit derivative at $R=1$ is$-6$, a unit, so $R$ is integral.
Finally
\[
 \kappa\equiv-6R-42R^7
             \equiv R-7(R-R^7)=R-7t\pmod{49}.
\]
\end{proof}

\begin{lemma}[The divided first row]\label{lem:first-row}
Put
\[
 E_i=(7^i-1)/6,\qquad
 S_r(v)=1+\sum_{i=1}^r(-1)^iv^{E_i}.
\]
The polynomial $V_1(1;v)$ is7-divisible and
\[
                       \overline{V_1(1;v)/7}=1+S_r(v). \tag{3}
\]
\end{lemma}
\begin{proof}
All series in this proof are formal in$t$, not evaluated functions.
Use the Artin--Hasse series
\[
 \mathcal E_7(X)=\exp\left(\sum_{i\ge0}X^{7^i}/7^i\right).
\]
Its integrality and inverse integrality are
\texttt{lem:artin-hasse-closure} of the
\href{ARITHMETIC_FRONTIER_TERMINAL_ORDERS.tex}{arithmetic note}.
Write $H=1-R$.
In the formal cyclic quotient $z^6=7v$, choose
$A_*(t)\in t\mathbb Z_7[v][[t]]$ uniquely by
\[
 A_*+vA_*^7+7^6v^8A_*^{49}+\cdots=H.
\]
Its linear coefficient equation has coefficient1, so recursive
formal inversion is integral. Then
$\exp(zH)=\mathcal E_7(zA_*)$ in this formal $t$-adic quotient.
At each fixed $t$-degree only finitely many $z$ powers occur.
Here $\mathcal C_{6,1}(\sum_jf_jz^j)=
\sum_{\ell\ge0}f_{6\ell+1}(7v)^\ell$, the coefficient of$z$
in the formal cyclic quotient.
Since $[X^7]\mathcal E_7=103/720\equiv2\pmod7$,
\[
 {\cal C}_{6,1}\exp(zH)
          \equiv A_*+14vA_*^7\pmod{49}.                 \tag{4}
\]
All higher section terms have a factor49.
There is no analytic boundary evaluation of either exponential.

Let $T=\sum_{i\ge0}t^{7^i}$, an integer series, and put
\[
 C(T)=\frac{T^7-\sum_{i\ge0}t^{7^{i+1}}}{7},\quad
 P(T)=T+3T^2+5T^3+5T^4+3T^5+T^6.
\]
These are integral and
$R=1+T+7R_1$, $R_1\equiv C(T)+P(T)\pmod7$.
Define integer polynomials $a_0=-1$,
$a_i=-1-va_{i-1}^7$ and $A_0=\sum_{i\ge0}a_it^{7^i}$.
With
$C(A_0)=(A_0^7-\sum_i a_i^7t^{7^{i+1}})/7$,
the equality $A_*+vA_*^7=H$ modulo49 gives
\[
 A_*=A_0+7A_1,\qquad
                   A_1\equiv-R_1-vC(A_0)\pmod7.
\]
The coefficient of $t^N$ in $(1+T)A_0$ is exactly zero:
$N=2L+1$ is not a sum of at most two powers of7.
By (2), the coefficient of $t^N$ in
$\kappa\,\mathcal C_{6,1}\exp(zH)$ is exactly $V_1(1;v)$.
Combining (4) and $\kappa=R-7t$ modulo49 therefore yields
\[
 \overline{V_1(1;v)/7}
 =[t^N]\{R_1(A_0-1-T)-v(1+T)C(A_0)
                    +2v(1+T)A_0^7-tA_0\}.             \tag{5}
\]
The carry series $C(T)$ and $C(A_0)$ have all exponents1
modulo6. Multiplication by a constant or a series supported
at powers of7 gives only exponents1 or2 modulo6,
whereas $N=3$ modulo6. Hence those carry terms vanish in (5).
After reduction modulo7, $A_0^7$ is supported on powers of7;
the last two terms likewise cannot contribute.
Only
\[
                     -5T^3+3T^2(A_0-T)
\]
survives from $P(T)(A_0-1-T)$.
Since three selected powers have no base-seven carry,
\[
 [t^N]T^3=3,\qquad [t^N]T^2A_0=a_0+2a_r=-1+2a_r.
\]
Thus (5) is $-27+6a_r=1-a_r$ in characteristic7.
The recursion gives $-a_r=S_r(v)$ modulo7.
This proves (3) uniformly in$r$.
\end{proof}

\begin{lemma}[The saturated boundary rows]\label{lem:rows}
For $\bar Q\in k[X]$ define
\[
 f(j)=(j+1)\bar Q(6-j),\qquad
 D_a=\sum_{j=0}^a(-1)^j\binom aj f(j).
\]
Then $\sum_{a=0}^5D_a=0$, $D_6=0$, and
$D_5=-[X^4]\bar Q$.
The rows $V_0,V_1$ are globally7-divisible. After dividing them
by7, the reduced six-row matrix on coefficient polynomials is
\[
 \begin{aligned}
 \bar V_a&=\binom a2(D_a-2SD_{a-1}+S^2D_{a-2}),
                                      &&2\le a\le5,\\
 R_0:=\overline{V_0/7}
      &=\bar Q(6)+v(2SD_5-S^2D_4),\\
 R_1:=\overline{V_1/7}
      &=S\bar Q(6)+\bar Q(5)+3vS^2[X^4]\bar Q ,
 \end{aligned}                                                   \tag{6}
\]
where $S=S_r(v)$.
\end{lemma}
\begin{proof}
For $0\le\ell\le N+4<7^{r+1}$ with $s_7(\ell)\le6$,
the top binomial argument in $G_Q(c-\ell)$ is
\[
 \tau-\ell/6,\qquad \tau=1+2\cdot7^{r+1}/6.
\]
Since $-1/6=\sum_{i\ge0}7^i$, its digits through position$r$
are the successive partial sums of the digits of$\ell$,
followed by adding1 to the units digit.
If $\ell_0\le5$, there is no carry: the top units digit is
$1+\ell_0$ and its$r$th digit is$s_7(\ell)$.
If $\ell_0=6$, the digit-sum bound forces$\ell=6$;
the top units digit is zero and Lucas' binomial vanishes.
As $N$ has units digit1 and$r$th digit2, Lucas' criterion gives
\[
 \overline{G_Q(c-\ell)}
          =(1+\ell_0)\binom{s_7(\ell)}2\bar Q(6-\ell_0)
          =\binom{s_7(\ell)}2f(\ell_0).
\]
In the backward difference for an index$j$ with digit sum$a\le6$,
only digitwise $\ell\le j$ survive modulo7.
Alternating differences of total higher order greater than two
kill this quadratic in the digit sum.
The low pattern contributes
$\binom a2(D_a-2D_{a-1}+D_{a-2})$.
A single higher digit1 contributes
$-(a-1)(D_{a-1}-D_{a-2})$.
Its normalized factorial weight is$a t_i$, where
$t_i=(-1)^iv^{E_i}$.
For two higher selections the difference is$D_{a-2}$.
Two selections at the same position have weight
$\binom a2t_i^2$; distinct positions have weight
$2\binom a2t_it_j$.
Thus their sum is $\binom a2(\sum_i t_i)^2D_{a-2}$.
Combining all three orders gives
\[
 \binom a2\{D_a-2(1+\sum_i t_i)D_{a-1}
                       +(1+\sum_i t_i)^2D_{a-2}\},
\]
which is the first line of (6), with $S=1+\sum_i t_i$.
It also gives $\bar V_1=0$; the constant row is zero as well.
The Newton expansion of $f$ at6 gives $\sum D_a=f(6)=0$,
and its degree gives the other assertions about $D_5,D_6$.

For the constant divided row, put
$\tau=(7N-1)/6=1+7(N-1)/6$.
The exact product identity gives
\[
 \frac{G_1(c)}7=(-1)^{N-2}\frac{\tau}{6N}
       \prod_{k=1}^{N-2}\left(1-\frac{7(N-1)}{6k}\right)
       \equiv1\pmod7 .
\]
Every correction has positive valuation.
The other terms of $V_0/7$ have digit sum exactly six.
Their sum is $v/6!$ times the formal row of digit sum six.
Using the first line of (6) at$a=6$ and $6!=-1$ in $k$
gives $v(2SD_5-S^2D_4)$.
This proves the second line.

Finally the exact cyclic raising recurrence, divided by7, gives
\[
 R_1(XQ)=(5+v\partial_v)R_1(Q)+R_0(Q).
\]
Starting with Lemma~\ref{lem:first-row} and using $vS'=S-1$
determines the row on $1,X,\ldots,X^4$:
\[
 [S+1,\ -S-2,\ S-3,\ -S-1,\ 3S^2v+S+2].
\]
This is precisely the last line of (6).
All divisions are by constant7 and preserve rank in characteristic
zero at every evaluated unit$v$.
\end{proof}

\begin{lemma}[A complete thirty-state certificate]\label{lem:orbit}
Put
\[
 g(T)=1+T+\cdots+T^5,\quad h(T)=g(T)+Tg'(T),
\]
\[
 q(T)=T^{10}+4T^9+3T^8+6T^7+5T^4+3T^3+6T^2+6T+2.
\]
In $R=\mathbb F_7[T]/(q(T))$, $T^5h(T)$ is a unit.
Set $v_*=-g'(T)/(T^5h(T))$ and
$x_0=1,\ x_{j+1}=1-v_*x_j^7$. Then
\[
 x_{30}=x_0,\qquad
                  \gcd(x_j-T,q(T))=1\quad(0\le j<30).    \tag{7}
\]
Thus $x_j-T$ is a unit for every nonnegative$j$.
\end{lemma}
\begin{proof}
This is a finite polynomial certificate; no range of terminal
indices is substituted for the unrestricted quantifier.
For explicit reproduction the residue representative is
\[
 v_*=6T^9+3T^8+6T^7+6T^6+4T^5+T^4+T^2+4T+5.
\]
Multiply by $T^5h$ modulo$q$ to obtain $-g'$.
The accompanying retained certificate lists all thirty $x_j$.
Its exact checker starts at1, applies the stated modular seventh
power recurrence thirty times, checks the return to1, and checks
each of the thirty Euclidean gcds in (7).
The return makes the whole sequence30-periodic by induction,
so the finite certificate covers every$j$, not just$j<30$.
Neither irreducibility of$q$ nor enumeration of field elements
is assumed.
\end{proof}

\begin{lemma}[No residual rank drop]\label{lem:rank}
The reduced matrix (6) has rank five at every unit
$v\in\bar{\mathbb F}_7$, for every$r\ge1$.
\end{lemma}
\begin{proof}
Suppose a nonzero coefficient polynomial is in its kernel.
If $S=0$, the first four rows force $D_2=\cdots=D_5=0$;
the constraint gives $D_1=-D_0$ and $R_0=D_0=0$,
a contradiction. Thus $S\ne0$.
The four second-order recurrence rows imply
\[
                         D_a=(A+Ba)S^a\quad(0\le a\le5).
\]
Put $w=vS^6\ne0$.
The hyperplane condition and the last two rows, after a unit
row scaling of the last row, are
\[
 \begin{pmatrix}
 g&Sg'\\1+w&-w\\1+S+Sw&-S+5Sw
 \end{pmatrix}\binom AB=0.
\]
Here $A\ne0$, and the second row gives $B/A=1+1/w$.
The first row becomes $hw+Sg'=0$.
If $h=0$, then $g=g'=0$, impossible since
$g=(S^6-1)/(S-1)$ has five distinct roots.
Hence $h\ne0$, $w=-Sg'/h$, and
$v=-g'/(S^5h)$.
The last row yields, after substitution,
\[
 -g'h(1+5S)+6S^2g'^2-h^2=0.
\]
Its monic scalar multiple is precisely $q(S)$.

There is consequently a map $R\to\bar{\mathbb F}_7$ sending
$T$ to this $S$ and $v_*$ to this $v$.
But $S_j(v)=1-vS_{j-1}(v)^7$, $S_0=1$.
Under this map $x_r-T$ therefore becomes
$S_r(v)-S=0$, contradicting its unit property in
Lemma~\ref{lem:orbit}. This proves rank five.
\end{proof}

\begin{theorem}[The second thin-digit family]\label{thm:family}
For every integer $r\ge1$, at terminal order $N=2\cdot7^r+1$,
the ordinary factorial moments of every nonzero complex
$X(A+BX^6)$ at orders $N-1,N$ cannot both vanish.
\end{theorem}
\begin{proof}
Let $T(t)=\sum_{k=0}^N\binom Nk(N)_{6k}t^k$.
The \href{PHYSICAL_PRIME_TOTAL_DEGREE_REDUCTION.tex}{physical-prime reduction},
\texttt{thm:radius}, confines every hypothetical multiple root
to $u=-1/t=7v$ with $v$ a unit.
The \href{gap_three_hurwitz_structure.tex}{all-gap critical-ideal identity},
\texttt{lem:critical-ideal}, supplies rank less than five for
the original cyclic matrix.
This implication transfers from $\mathbb C$ to $\mathbb C_7$
algebraically: all coefficients and cyclic minors are rational
at fixed$N$, every multiple root is algebraic, and the complex
minor vanishing transports through its number-field embeddings.
No exponential is evaluated7-adically.
The monic triangular column change to
$\Pi_N(A)(1,A,\ldots,A^4)1$ and the row normalizations (1)
therefore give a nonzero integral primitive kernel vector.
Dividing rows0 and1 by7 preserves this relation.
Its reduction contradicts Lemma~\ref{lem:rank}.

Thus $T$ is squarefree.
The normalized consecutive moments are
$T+(6/N)tT'$ and $T-(1/N)tT'$;
their coefficient determinant is $-7/N\ne0$ and $T(0)=1$.
They cannot have a common root.
Nonzero degenerate monomials have no vanishing positive moment.
\end{proof}
This is an unbounded terminal family.
It does not prove ordinary gap six at every index,
degree-seven binomial SFC, all-binomial SFC, or unrestricted FC/SFC.
The written all-digit argument and the complete finite orbit,
rather than a terminal-index prefix, supply its quantifiers.
\end{document}
